\section{Implications for Model Evaluation}
\label{sec:framework}

The experiments do not support a rule that assigns one model family to a
retail dataset from a few summary statistics. They do support a relatively
simple evaluation procedure. Its purpose is to prevent differences in samples,
metrics, and baseline effort from being mistaken for differences between
models.

\subsection{A common-panel comparison}
\label{subsec:validation-checklist}

The first step is to choose an accuracy measure that represents the decision
being supported. Aggregate or volume-weighted errors may suit warehouse and
category planning, while item-level replenishment may require per-series or
service-level measures. Safety-stock decisions should be evaluated with
quantile loss or the relevant inventory cost rather than point-forecast error
alone. On intermittent series, both the number of zero-demand periods and the
denominator of percentage measures should be inspected.

The candidate models should then be run on identical series, origins,
horizons, and actual values. Seasonal Naive provides a minimum reference. A
global LightGBM model with available covariates represents a practical
supervised baseline, and at least one alternative multi-step strategy should
be checked when the horizon is long. One or two zero-shot foundation models
can be added without an extensive search over checkpoints. Saving individual
forecasts permits a paired bootstrap comparison and makes later error analysis
possible.

The final choice should consider the size and uncertainty of the accuracy
difference rather than the rank alone. If the paired interval is practically
indistinguishable from zero, implementation effort, latency, licensing, and
maintenance become reasonable tie-breakers. Comparisons based on different
series counts or forecast windows should be labelled as descriptive and should
not determine deployment.

\subsection{Factors that change the comparison}
\label{subsec:framework-moderators}

Metric choice produced the most consistent difference in the public results.
The strong pattern for Scaled Quantile Loss does not imply an equally strong
advantage in \WAPE{} or \MASE{}. A team should therefore avoid an overall
benchmark rank that averages measures with different operational meanings.

Demand intermittency can affect both model behaviour and metric aggregation.
Our direct--recursive LightGBM results differ between samples and between
per-series and aggregate \WAPE{}. The zero fraction is consequently useful for
stratifying an evaluation panel, but the three datasets do not establish a
threshold at which one strategy should replace another.

Covariates create a further asymmetry. Prices, promotions, holidays, weather,
and stock availability can give a supervised model information that a
univariate foundation model does not receive. A zero-shot model may still be
valuable when such features are unavailable or costly to maintain, but this is
an engineering trade-off rather than an accuracy advantage under equal
information.

Finally, the measured laptop and Azure runtimes should not be used as a general
cost comparison. Throughput depends on batching, model implementation,
hardware, and the number and length of series. These quantities need to be
measured on the intended workload. Licence restrictions also matter:
Moirai~2.0-Small, for example, is distributed under a non-commercial licence.

\subsection{Deployment choice}
\label{subsec:framework-decision-rule}

A foundation model is a reasonable choice when it matches or improves the
target metric on the common panel, its inference requirements are acceptable,
and avoiding a feature and training pipeline has material value. A supervised
baseline is preferable when it performs better, already operates reliably in
production, or benefits from important covariates that the candidate
foundation model cannot use. If neither improves on Seasonal Naive, the first
response should be to inspect the data and evaluation design rather than to
increase model complexity.
